\section{Описание предложенного метода}
\label{sec:method}

Теория из раздела~\ref{sec:notation} даёт математическую формулировку задачи и доказывает, что искомое поле скоростей~$u_t$ можно обучить симуляционно-свободно, а при линейном интерполянте одного шага метода Эйлера достаточно для численного интегрирования. Настоящий раздел переводит эти результаты в практический алгоритм: фиксируется конкретный класс параметризаций~$f_\theta$ (архитектура TimeResNet), описывается схема обучения и инференса, и формулируются базовые методы сравнения.

\subsection{Общая схема метода MFM}

Метод \textbf{Medical Flow Matching~(MFM)} объединяет три ключевых компонента: (1)~предобработку и регистрацию парных КТ-данных, (2)~обучение нейронной сети~$f_\theta$ предсказывать поле скоростей в рамках условного флоу матчинга и (3)~инференс через однократное интегрирование ОДУ.

Схема работы метода представлена на рисунке~\ref{fig:flow}. Обучение выполняется следующим образом: для каждой пары $(\mathbf{x}_0, \mathbf{x}_1)$ и случайного времени $t \sim \mathcal{U}[0,1]$ формируется промежуточный образец $\mathbf{x}_t = (1-t)\mathbf{x}_0 + t\mathbf{x}_1$ и целевая скорость $\mathbf{v}_t = \mathbf{x}_1 - \mathbf{x}_0$; сеть $f_\theta$ минимизирует $\ell_2$-расстояние между предсказанной и целевой скоростями. При инференсе интегрирование уравнения~(\ref{eq:ode_flow}) от $t=0$ до $t=1$ даёт синтетический контрастный срез; от $t=1$ до $t=0$~--- нативный. Единственная сеть обслуживает оба направления, что вдвое сокращает затраты на обучение и хранение весов.

\begin{figure}[htbp]
    \centering
    \includegraphics[width=0.85\linewidth]{flow.pdf}
    \caption{Схема обучения и инференса MFM. Сверху: линейная интерполяция между нативным и контрастным срезами при разных значениях $t$. Снизу: обучение сети предсказывать постоянную скорость $\mathbf{v} = \mathbf{x}_1 - \mathbf{x}_0$.}
    \label{fig:flow}
\end{figure}

\subsection{Условный флоу матчинг}
\label{sec:fm}

Напомним формальную постановку, опираясь на обозначения раздела~\ref{sec:notation}. Пусть $\pi_0 = \pi_N$ (нативное) и $\pi_1 = \pi_C$ (контрастное) при трансляции в направлении нативное~$\to$~контрастное; роли обмениваются для обратного направления. Линейный интерполянт~(\ref{eq:linear_interp}) определяет вероятностный путь $\{\pi_t\}$, а потеря условного флоу матчинга~(\ref{eq:cfm_loss}) обучает сеть предсказывать постоянную целевую скорость~$\mathbf{x}_1 - \mathbf{x}_0$.

Принципиальное отличие от диффузионных моделей состоит в том, что:
\begin{itemize}
    \item целевая скорость $\mathbf{v}_t = \mathbf{x}_1 - \mathbf{x}_0$ не зависит от $t$~--- сети достаточно выучить одну функцию, а не $T$ различных;
    \item потеря~(\ref{eq:cfm_loss}) вычисляется аналитически, без численного интегрирования ОДУ;
    \item при инференсе достаточно одного шага метода Эйлера, поскольку выученное поле скоростей практически линейно.
\end{itemize}

\paragraph{Двунаправленная трансляция.} Так как линейный интерполянт~(\ref{eq:linear_interp}) симметричен по $\mathbf{x}_0$ и $\mathbf{x}_1$, трансляция в обоих направлениях обеспечивается одной сетью простой заменой пределов интегрирования: интегрирование от 0 к 1 выполняет нативное~$\to$~контрастное, от 1 к 0~--- контрастное~$\to$~нативное. В процессе обучения пары предъявляются в обоих порядках, что дополнительно обогащает разнообразие обучающих примеров.

\subsection{Архитектура TimeResNet}
\label{sec:arch}

TimeResNet представляет собой U-Net энкодер--декодер с двумя специализированными модификациями, разработанными применительно к задаче условного флоу матчинга (рисунок~\ref{fig:arch}).

\begin{figure}[htbp]
    \centering
    \includegraphics[width=\linewidth]{timeresnet.pdf}
    \caption{Архитектура TimeResNet: U-Net с поблочной аффинной временной модуляцией и модулем самовнимания на боттлнеке. Единственная сеть используется для обоих направлений трансляции.}
    \label{fig:arch}
\end{figure}

\subsubsection{(а) Поблочное внедрение временного сигнала}

Скалярный параметр $t \in [0,1]$ кодируется посредством синусоидального позиционного кодирования~\cite{vaswani2017attention} и преобразуется двухслойным персептроном с активацией SiLU (рисунок~\ref{fig:embedding}). На выходе MLP получаются векторы $\mathbf{A}, \mathbf{B} \in \mathbb{R}^C$, специфичные для данного ResNet-блока. Скрытая активация $\mathbf{h}$ после первой пары GroupNorm--SiLU--Conv модулируется аффинно:
\begin{equation}
    \mathbf{h} \leftarrow \mathbf{A} \odot \mathbf{h} + \mathbf{B},
    \label{eq:time_mod}
\end{equation}
где $\odot$~--- поэлементное умножение с трансляцией по пространственным измерениям.

\begin{figure}[htbp]
    \centering
    \includegraphics[height=0.7\textheight,keepaspectratio]{Embeding.pdf}
    \caption{Модуль временного эмбеддинга: синусоидальное кодирование $\to$ двухслойный MLP $\to$ аффинная модуляция (масштаб~$\mathbf{A}$, сдвиг~$\mathbf{B}$; формула~(\ref{eq:time_mod})).}
    \label{fig:embedding}
\end{figure}

Внедрение временного сигнала \emph{в каждый блок} принципиально отличает TimeResNet от наивного подхода, при котором $t$ конкатенируется лишь с входом сети. Следующее утверждение объясняет, почему это различие существенно с точки зрения градиентного сигнала.

\begin{proposition}[Сохранение временн\'{о}го сигнала на всех глубинах]
\label{prop:time_signal}
Обозначим скрытые активации в блоках $l=1,\ldots,L$ через $\mathbf{h}_l$, а потерю CFM~--- через $\mathcal{L}$. Пусть~$\mathcal{M}_l(\mathbf{h}_l;t)$~--- оператор аффинной временной модуляции в блоке~$l$, определённый формулой~\eqref{eq:time_mod}, с обучаемой проекцией $W_l$, не сводящейся к нулевой. Тогда:
\begin{enumerate}
    \item[\rm (i)] При поблочном внедрении на каждой глубине~$l$ существует ненулевая компонента градиента по~$t$:
    \begin{equation}
        \frac{\partial \mathcal{L}}{\partial t}
        = \sum_{l=1}^{L} \frac{\partial \mathcal{L}}{\partial \mathbf{h}_l} \cdot \frac{\partial \mathcal{M}_l(\mathbf{h}_l; t)}{\partial t} \neq 0
        \qquad \text{в общем положении.}
        \label{eq:dLdt_perblock}
    \end{equation}
    \item[\rm (ii)] При наивной конкатенации $t$ только со входом сети ($\mathbf{h}_0=[\mathbf{x};t]$) выполнено
    \begin{equation}
        \frac{\partial \mathbf{h}_l}{\partial t} = 0 \qquad \forall\, l>1,
        \label{eq:dh_concat}
    \end{equation}
    т.\,е.\ временной сигнал переносится через сеть лишь через скрытые признаки $\mathbf{h}_1,\mathbf{h}_2,\ldots$ опосредованно, без явного канала, и легко затухает по глубине.
\end{enumerate}
\end{proposition}

\begin{proof}
\textbf{(i)} Применим цепное правило к функционалу $\mathcal{L}(\theta;t)$. Поскольку оператор $\mathcal{M}_l$ зависит от $t$ через коэффициенты $\mathbf{A}_l(t),\mathbf{B}_l(t)$, прямой подсчёт даёт
\[
    \frac{\partial \mathcal{M}_l(\mathbf{h}_l;t)}{\partial t}
    = \frac{\partial \mathbf{A}_l}{\partial t}\odot\mathbf{h}_l + \frac{\partial \mathbf{B}_l}{\partial t}, \qquad
    \frac{\partial (\mathbf{A}_l,\mathbf{B}_l)}{\partial t} = W_l\cdot\mathrm{MLP}'(\gamma(t))\cdot\gamma'(t).
\]
Синусоидальное кодирование~$\gamma$ обладает почти всюду ненулевой производной $\gamma'(t)$, а $W_l$ по предположению не сводится к нулю. Поэтому $\partial\mathcal{M}_l/\partial t\neq 0$ почти всюду. Просуммировав по~$l=1,\ldots,L$ вклады по цепному правилу, получаем~\eqref{eq:dLdt_perblock}.

\textbf{(ii)} При наивной конкатенации сигнал $t$ входит в сеть лишь через $\mathbf{h}_0$. Для последующих блоков $\mathbf{h}_l = g_l(\mathbf{h}_{l-1})$ не зависит от $t$ напрямую, а только через композицию $\mathbf{h}_l = g_l\circ\cdots\circ g_1(\mathbf{h}_0(t))$. Тогда
\[
    \frac{\partial \mathbf{h}_l}{\partial t} = \prod_{k=1}^{l} J_{g_k}\,\cdot\,\frac{\partial \mathbf{h}_0}{\partial t},
\]
где $J_{g_k}$~--- якобиан блока. Произведение норм $\prod_k\lVert J_{g_k}\rVert$ может экспоненциально затухать с глубиной (классическая проблема исчезающего градиента). В отличие от~\eqref{eq:dLdt_perblock}, здесь не существует независимого ``канала'' для $t$ в глубоких блоках. Утверждение~\eqref{eq:dh_concat} выполнено в смысле отсутствия прямой зависимости.
\end{proof}

Абляционный эксперимент (раздел~\ref{sec:ablation}) количественно подтверждает следствие утверждения~\ref{prop:time_signal}: замена поблочного внедрения входовой конкатенацией приводит к росту MAE на 1.77\,HU.

\subsubsection{(б) Свёрточное самовнимание на боттлнеке}

Между двумя боттлнековыми ResNet-блоками вставляется модуль самовнимания~\cite{yang-etal-2019-convolutional} (рисунок~\ref{fig:self_attn}). Обозначим вход модуля как $\mathbf{X}\in\mathbb{R}^{C\times H\times W}$, где $C$~--- число каналов, $H,W$~--- пространственные размеры боттлнека.

\paragraph{Шаг 1: нормализация.}
Применяется групповая нормировка~\cite{wu2018group} с $G=32$ группами по каналам:
\begin{equation}
    \tilde{\mathbf{X}} = \mathrm{GroupNorm}_{G}(\mathbf{X}), \qquad
    \tilde{\mathbf{X}}_{c,h,w} = \gamma_c\,\frac{\mathbf{X}_{c,h,w}-\mu_g}{\sqrt{\sigma_g^2+\varepsilon}} + \beta_c,
    \label{eq:groupnorm}
\end{equation}
где $g=\lfloor c\,G/C\rfloor$~--- номер группы, $\mu_g,\sigma_g^2$~--- среднее и дисперсия активаций по каналам группы~$g$ и всем пространственным позициям, $\gamma_c,\beta_c$~--- обучаемые масштаб и сдвиг.

\paragraph{Шаг 2: проекции запросов, ключей и значений.}
Тензор $\tilde{\mathbf{X}}$ преобразуется тремя независимыми свёртками $1\times1$ в тензоры $\mathbf{Q},\mathbf{K},\mathbf{V}\in\mathbb{R}^{C\times H\times W}$ с обучаемыми ядрами $\mathbf{W}^Q,\mathbf{W}^K,\mathbf{W}^V\in\mathbb{R}^{C\times C}$:
\begin{equation}
    \mathbf{Q} = \mathbf{W}^Q\!\ast\tilde{\mathbf{X}}, \qquad
    \mathbf{K} = \mathbf{W}^K\!\ast\tilde{\mathbf{X}}, \qquad
    \mathbf{V} = \mathbf{W}^V\!\ast\tilde{\mathbf{X}},
    \label{eq:qkv_proj}
\end{equation}
где $\ast$~--- операция свёртки. Свёртка $1\times1$ эквивалентна аффинному преобразованию вектора каналов в каждой пространственной позиции независимо.

\paragraph{Шаг 3: развёртывание в последовательность токенов.}
Пространственные оси объединяются в одну: $N := HW$. Каждая позиция трактуется как отдельный токен, и тензоры приводятся к матрицам $Q,K,V\in\mathbb{R}^{N\times C}$:
\begin{equation}
    Q_{n,:} = \mathbf{Q}_{:,h(n),w(n)}, \qquad
    K_{n,:} = \mathbf{K}_{:,h(n),w(n)}, \qquad
    V_{n,:} = \mathbf{V}_{:,h(n),w(n)},
    \label{eq:flatten}
\end{equation}
где $(h(n),w(n))$~--- линейная развёртка пространственной решётки, $n=1,\ldots,N$.

\paragraph{Шаг 4: масштабированное скалярное произведение.}
Вычисляется матрица попарных сходств токенов и применяется softmax по строкам:
\begin{equation}
    \mathbf{S} = \frac{Q K^\top}{\sqrt{C}} \in\mathbb{R}^{N\times N}, \qquad
    \mathbf{A}_{n,m} = \frac{\exp(\mathbf{S}_{n,m})}{\sum_{m'=1}^{N}\exp(\mathbf{S}_{n,m'})},
    \label{eq:attn_weights}
\end{equation}
\begin{equation}
    \mathrm{Attention}(Q,K,V) = \mathbf{A}\,V \in\mathbb{R}^{N\times C}.
    \label{eq:attn_out}
\end{equation}
Деление на $\sqrt{C}$~--- классический трюк~\cite{vaswani2017attention}: при больших~$C$ скалярное произведение $Q_n K_m^\top$ имеет дисперсию порядка~$C$, и без масштабирования аргумент softmax насыщается, обнуляя градиент. Поэлементная стоимость~\eqref{eq:attn_weights}--\eqref{eq:attn_out} равна $\mathcal{O}(N^2 C) = \mathcal{O}(H^2 W^2 C)$~--- квадратична по числу токенов, но применяется только на боттлнеке, где $H,W$ малы.

\paragraph{Шаг 5: сворачивание и остаточное сложение.}
Результат~\eqref{eq:attn_out} сворачивается обратно в форму~$C\times H\times W$ обратной развёрткой~\eqref{eq:flatten}, проецируется выходной свёрткой $1\times1$ с ядром $\mathbf{W}^O\in\mathbb{R}^{C\times C}$ и прибавляется к исходному входу:
\begin{equation}
    \mathbf{Y} = \mathbf{X} + \mathbf{W}^O\!\ast\mathrm{Reshape}_{C\times H\times W}\!\bigl(\mathrm{Attention}(Q,K,V)\bigr).
    \label{eq:attn_residual}
\end{equation}
Остаточное соединение обеспечивает тождественную инициализацию модуля: при $\mathbf{W}^O=0$ модуль не меняет тензор, что облегчает оптимизацию на ранних эпохах~\cite{he2016deep}.

\begin{figure}[htbp]
    \centering
    \includegraphics[width=0.7\linewidth]{self_attention.pdf}
    \caption{Модуль свёрточного самовнимания на боттлнеке: GroupNorm $\to$ $1\times1$~Conv ($Q, K, V$) $\to$ развёртывание $\to$ масштабированное скалярное произведение~\eqref{eq:attn_weights}--\eqref{eq:attn_out} $\to$ остаточное сложение~\eqref{eq:attn_residual}.}
    \label{fig:self_attn}
\end{figure}

Внимание применяется на \emph{наименьшем} пространственном разрешении энкодера, что делает квадратичную по~$N=HW$ стоимость~\eqref{eq:attn_weights} разумной. Уточним этот выигрыш формально.

\begin{proposition}[Вычислительная сложность самовнимания на боттлнеке]
\label{prop:attn_complexity}
Пусть исходное разрешение входа равно $H\times W$, глубина энкодера~--- $k$ (количество понижающих этапов с коэффициентом~2 на этап), а на каждом этапе уменьшение происходит независимо по высоте и ширине. Тогда число токенов в боттлнеке
\begin{equation}
    N_{\mathrm{bottle}} = \frac{HW}{4^k},
    \label{eq:Nbottle}
\end{equation}
а отношение стоимостей самовнимания на боттлнеке и на полном разрешении составляет
\begin{equation}
    \frac{\mathrm{Cost}_{\mathrm{bottle}}}{\mathrm{Cost}_{\mathrm{full}}}
    = \frac{\mathcal{O}(N_{\mathrm{bottle}}^2\,C)}{\mathcal{O}((HW)^2\,C)}
    = \frac{1}{4^{2k}}.
    \label{eq:attn_ratio}
\end{equation}
В частности, при~$k=4$ боттлнековое самовнимание примерно в~$4^{8}=65\,536$ раз дешевле полноразмерного.
\end{proposition}

\begin{proof}
Каждый из~$k$ этапов энкодера уменьшает пространственное разрешение в~2 раза по обеим осям, что соответствует уменьшению числа пикселей в~$4=2^2$ раза. После~$k$ этапов имеется $N_{\mathrm{bottle}}=HW/4^k$ токенов~--- получено~\eqref{eq:Nbottle}.

Стоимость стандартного самовнимания на множестве из~$N$ токенов и размерности признака~$C$ доминируется матричным произведением $QK^\top\in\mathbb{R}^{N\times N}$ и его умножением на $V\in\mathbb{R}^{N\times C}$, что в сумме даёт $\mathcal{O}(N^2 C)$ операций. Подставив $N=HW$ для полноразмерного варианта и $N=N_{\mathrm{bottle}}$ для боттлнекового и поделив, получаем~\eqref{eq:attn_ratio}.

Численный пример: $k=4 \Rightarrow$ $4^{2k}=4^{8}=65\,536$, что и заявлено.
\end{proof}

В отличие от оконного внимания Swin~\cite{hatamizadeh2021swin}, предложенный модуль не ограничен фиксированным размером окна и способен захватывать долгодистанционные зависимости во всей карте признаков боттлнека~--- именно это свойство критично при трансляции КТ, где удалённые анатомические структуры (например, сосудистое русло) скоррелированы по контрастности. Удаление этого модуля, согласно абляционному исследованию (раздел~\ref{sec:ablation}), ведёт к росту MAE на 1.71\,HU.

\subsection{Базовые методы сравнения}
\label{sec:baselines}

Для всестороннего сравнения рассматриваются три семейства методов.

\paragraph{Регрессионные сети.} Прямолинейное обучение сети $g_\theta$ минимизировать $\ell_1$-потерю:
\[
    \mathcal{L}_{\mathrm{reg}}(\theta) = \mathbb{E}\bigl[\lVert g_\theta(\mathbf{x}_C) - \mathbf{x}_N \rVert_1\bigr].
\]
Тестируются четыре архитектуры: UNet~\cite{ronneberger2015u}, SegResNet~\cite{myronenko20183d}, SwinUNETR~\cite{hatamizadeh2021swin}, UMambaBot~\cite{ma2024u}.

\paragraph{DDPM.} Условная диффузионная модель, конкатенирующая исходный срез с зашумлённым целевым и использующая DDIM~\cite{song2020denoising} для ускоренного семплирования. Параметры выбраны так, чтобы обеспечить наилучшее качество с разумным временем инференса.

\paragraph{SMILE.} Метод SMILE~\cite{liu2025see}~--- метод анатомически-управляемого диффузионного контрастного усиления, оцениваемый с использованием опубликованных авторами весов и рекомендованных настроек.
